% !TeX root = ../papers/nc-p3.tex
\section{四次K3曲面上の切断と二つの楕円鉛筆}
\label{sec:finite-quartic-k3-cut}
\begin{theorem}\label{thm:finite-quartic-k3-cut}
定理\ref{thm:finite-quartic-grassmann-windows}の両局所自由部分で、
一般の二平面 $P\subset H$ に対し $T=\{x:f(x)\cap P\ne0\}$ は滑らかな四次K3曲面である。
$h$ を超平面類とする。$T$ 上には二つの基点のない楕円鉛筆 $|C|,|F|$ があり、
それぞれの一般ファイバーから元の窓とその二移動を復元できる。
\[
 C+F=4h,\quad hC=hF=8,\quad C^2=F^2=0,\quad CF=32 .
\]
$D=2h-C=F-2h$ と置く。導来圏 $D^b(T)$ の四対象
\[
 G=(\mathcal O,\mathcal O(h),\mathcal O(D)[1],\mathcal O(h+D)[1])
\]
の次数零自己準同型代数は、積の向きを経路順に読んだとき元の小さい四層切断 $A$ である。
次数つき自己準同型代数全体は二次の双対を加えた $A\ltimes DA[-2]$ で、
各双対辺に補経路長を与えれば元の有限Frobenius代数を得る。
\end{theorem}
\begin{proof}
一般切断 $s\in H$ は滑らかな楕円八次曲線を零点に持つ。
$P=\langle s,t\rangle$ とすれば $T$ の方程式は $s\wedge t$ である。
生成性により $s$ の零曲線の外でこの線型系は自由であり、そこで一般の成員は滑らかである。
零曲線上で $t$ も零となる条件の余次元は固定点で2なので、一般の $t$ はその全点で非零である。
$s$ の微分は零曲線上で束へ全射だから $s\wedge t$ の微分もそこで消えない。
こうして得た $T$ の全点で $f(x)\cap P$ は直線であり、これと $f(x)$ の $H/P$ における像直線が
二つの射 $T\to\mathbb P^1$ を与える。前者の繊維は $P$ の指定切断の零曲線である。
後者も $P^\perp$ の指定切断の零曲線であり、双対評価列に同じ議論を適用できる。
一般の二平面を取っているので両側とも一般繊維は滑らかで連結である。
核二平面束の行列式を二直線から計算すれば $C+F=4h$。二鉛筆の自己交叉と曲線の次数から残りの数値が従う。

$k=1,2,3$ に対し $H^1(T,kh-C)=H^1(\mathbb P^3,I_C(k))$ が乗法と両立する。
これは超曲面のイデアル列と射影三空間の直線束の中間消滅による。
また $H^0(T,kh-C)=0$ は $C$ が三次以下の曲面に含まれないことによる。
$C-kh=(4-k)h-F$ に同じ事実を用いれば $H^2(T,kh-C)=0$ である。
従って二つの隣接する無移動直線束の間の順方向射は $H^0(h)$、
シフトをまたぐ順方向射は上の三つのRao空間で、全て導来次数零にある。
合成は二端での $H^0(h)$ の作用、すなわち三段加群の可換乗法である。
一次関係と二回乗法の復元によりこれが $A$ を与える。
逆方向はK3のSerre双対で対応する順方向の双対が次数2にあるだけである。
対角には恒等射とその次数2の双対がある。trace対合の合成との両立性と次数4の消滅から
全積が上の自明拡大となる。この議論に高次積の消滅は必要ない。
\end{proof}

\begin{proposition}\label{prop:quartic-k3-two-pencil-twists}
$T_M$ を球面対象 $M$ に沿うtwist、$L_a$ を $\mathcal O_T(a)$ とのテンソル積とする。
積の右側が先に作用する。$A_0=T_{\mathcal O}T_{\mathcal O(h)}$ とすると
\[
 \Phi=T_{G_0}T_{G_1}T_{G_2}T_{G_3},\qquad
 \Phi[-2]= A_0 L_{-C} A_0^{-1} L_{-F}.
\]
数値的には $v=(0,F,-4)$ を $\Phi$ が固定し、$A_0^{-1}$ はこれを $(0,-C,-12)$ に送る。
\end{proposition}
\begin{proof}
$B_0=T_{\mathcal O(-2h)}T_{\mathcal O(-h)}$ とする。射影三空間の連続した四直線束を四次K3へ制限すると
その順序のtwistの積は $L_{-4h}[2]$ となることを確認する。
各twistの核は $M^*\boxtimes M\to\mathcal O_\Delta$ の錐である。
積を展開すると、恒等項の手前は四直線束の真に増加する列から作る評価複体である。
長さ1から3の差の順方向の導来Homは、射影三空間から制限して同型である
（超曲面列の負の捩りは $-3,-2,-1$）。よって評価複体は射影三空間での同じ複体の導来引戻しである。
射影三空間ではこの複体が対角層を与える。実際、四直線束は充満例外列なので
順方向の射を用いた正規化バー分解の評価が恒等関手の核を解消する
（同じことは右端から各例外成分を順に除く射影の積が零になることからも分かる）。
従って曲面での積は引戻した対角層から曲面自身の対角層への増大の錐である。
四次超曲面の一項のKoszul分解から、この増大は次数零で同型、
唯一の負次数コホモロジーは次数 $-1$ の対角上の $\mathcal O_T(-4h)$ である。
よって錐は $L_{-4h}[2]$ を与え、特に $B_0 A_0=L_{-4h}[2]$。
twistは球面対象のシフトに依存せず、テンソル積による共役と両立する。
$D=F-2h$ より $\Phi=A_0 L_F B_0 L_{-F}$、従って $F-4h=-C$ を代入して表示式を得る。
Mukai対合を $(r,c,s)\cdot(r',c',s')=cc'-rs'-r's$ とすれば
twistとその逆はともに $x\mapsto x+(x\cdot v(M))v(M)$ と作用する。
$v$ にまず $\mathcal O$、次に $\mathcal O(h)$ の反射を施せば $(4,F,0)$ を経て $(0,-C,-12)$ となる。
$F^2=C^2=0$ と分解式で固定性が従う。
\end{proof}
